\begin{name}
	{NGUYÊN HÀM - TÍCH PHÂN}
	{KT TÍCH PHÂN}
	{\tentruong}
	{\thoigian}
\end{name}
\setcounter{ex}{0}\setcounter{bt}{0}
\Opensolutionfile{ans}[ans/ans-2-B12-De2-NLC]
\TN
\begin{ex}%[Cau-1]%[2D4N2-1]
	Cho hàm số $y=f(x)$ liên tục trên khoảng $K$ và $a$, $b$, $c\in K$. Mệnh đề nào sau đây \textbf{sai}?
	\choice
	{$\displaystyle \int\limits_{a}^{a} f(x)\mathrm{\,d}x=0$}
	{$\displaystyle \int\limits_{a}^{b} f(x)\mathrm{\,d}x=\displaystyle \int\limits_{a}^{b} f(t) \mathrm{\,d}t$}
	{$\displaystyle \int\limits_{a}^{b} f(x)\mathrm{\,d}x=-\displaystyle \int\limits_{b}^{a} f(x) \mathrm{\,d}x$}
	{\True $\displaystyle \int\limits_{a}^{b} f(x) \mathrm{\,d}x + \displaystyle \int\limits_{c}^{b} f(x) \mathrm{\,d}x=\displaystyle \int\limits_{a}^{c} f(x) \mathrm{\,d}x$}
	\loigiai{
		Mệnh đề sai là $\displaystyle \int\limits_{a}^{b} f(x)\mathrm{\,d}x + \displaystyle \int\limits_{c}^{b} f(x) \mathrm{\,d}x=\displaystyle \int\limits_{a}^{c} f(x) \mathrm{\,d}x$.
	}
\end{ex}
\begin{ex}%[Cau-2]%[2D4N2-1]
	Cho hàm số $f(x)$ liên tục trên $\mathbb{R}$ và $F(x)$ là nguyên hàm của $f(x)$, biết $\displaystyle \int\limits_{0}^{9} f(x)\mathrm{\,d}x=9$ và $F(0)=3$. Tính $F(9)$.
	\choice
	{$F(9)=-6$}
	{$F(9)=6$}
	{\True $F(9)=12$}
	{$F(9)=-12$}
	\loigiai{
		Ta có $I=\displaystyle\int\limits_{0}^{9} f(x)\mathrm{\,d}x = F(x)\Big|_0^9 = F(9)- F(0)=9\Leftrightarrow F(9)=9 + F(0)=9 + 3= 12$.
	}
\end{ex}
\begin{ex}%[Cau-3]%[2D4N2-2]
	Tính tích phân $I=\displaystyle\int\limits_{0}^{1} x^{2018}(1 + x)\mathrm{\,d}x$?
	\choice
	{$I=\dfrac{1}{2017}+\dfrac{1}{2018}$}
	{$I=\dfrac{1}{2018}+\dfrac{1}{2019}$}
	{$I=\dfrac{1}{2020}+\dfrac{1}{2021}$}
	{\True $I=\dfrac{1}{2019} + \dfrac{1}{2020}$}
	\loigiai{
		Ta có $I=\displaystyle \int\limits_{0}^{1} x^{2018}(1 + x) \mathrm{\,d}x= \displaystyle \int\limits_{0}^{1} \left(x^{2018} + x^{2019}\right) \mathrm{d}x= \left(\dfrac{x^{2019}}{2019} + \dfrac{x^{2020}}{2020}\right)\Bigg|_0^1 = \dfrac{1}{2019}+\dfrac{1}{2020}$.
	}
\end{ex}
\begin{ex}%[Cau-4]%[2D4N2-4]
	Tính $\displaystyle \int\limits_{0}^{1} 2\mathrm{e}^{x} \mathrm{\,d}x$?
	\choice
	{$I=\mathrm{e}^2-2\mathrm{e}$}
	{$I=2\mathrm{e}$}
	{$I=2\mathrm{e}+2$}
	{\True $I=2\mathrm{e}-2$}
	\loigiai{
		Ta có $\displaystyle \int\limits_{0}^{1} 2\mathrm{e}^{x} \mathrm{\,d}x = 2\mathrm{e}^x\Big|_0^1=2\mathrm{e}-2$.
	}
\end{ex}
\begin{ex}%[Cau-5]%[2D4N2-3]
	Cho $a\in \left(0;\dfrac{\pi}{2}\right)$. Tính $J=\displaystyle \int\limits_{0}^{a} \dfrac{29}{\cos^2 x} \mathrm{\,d}x$ theo $a$.
	\choice
	{$J=-29\tan a$}
	{$J=\dfrac{1}{29}\tan a$}
	{$J=29\cot a$}
	{\True $J=29\tan a$}
	\loigiai{
		Ta có $J=\displaystyle \int\limits_{0}^{a} \dfrac{29}{\cos^2 x}\mathrm{\,d}x=29\tan x\Big|_0^a = 29\tan a - 29\tan 0 = 29\tan a$.
	}
\end{ex}
\begin{ex}%[Cau-6]%[2D4H2-5]
	Tích phân $I=\displaystyle \int\limits_{-1}^{2}\left|x^2 - 2x\right|\mathrm{d}x$ có giá trị là
	\choice
	{\True $I=\dfrac{8}{3}$}
	{$I=\dfrac{4}{3}$}
	{$I=0$}
	{$I=-\dfrac{4}{3}$}
	\loigiai{
		Ta có $x^2 - 2x=0 \Leftrightarrow x=0$ hoặc $x=2$.\\
		Bảng xét dấu
		\begin{center}
			\begin{tikzpicture}
				\tkzTabInit[nocadre=false,lgt=2.5,espcl=2.5,deltacl=0.6]
				{$x$/0.7,$x^2-2x$/0.7}
				{$-\infty$,$0$,$2$,$+\infty$}
				\tkzTabLine{,+,0,-,0,+,}   
			\end{tikzpicture}
		\end{center}
		\begin{align*}
			I & = \displaystyle \int\limits_{-1}^{2} \left|x^2 - 2x\right| \mathrm{d}x = \displaystyle \int\limits_{-1}^{0} \left(x^2 - 2x\right) \mathrm{d}x - \displaystyle \int\limits_{0}^{2} \left(x^2 - 2x\right) \mathrm{d}x\\
			& = \left(\dfrac{x^3}{3} - x^2\right)\bigg|_{-1}^0 - \left(\dfrac{x^3}{3} - x^2\right)\bigg|_{0}^2 \\
			& = \left[0 - \left(-\dfrac{1}{3}-1\right)\right] - \left[\left(\dfrac{8}{3}-4\right)-0\right]\\
			& = \dfrac{8}{3}.
		\end{align*}
	}
\end{ex}
\begin{ex}%[Cau-7]%[2D4N2-3]
	Tính tích phân $I=\displaystyle \int\limits_{0}^{\tfrac{\pi}{4}} \sin x \mathrm{\,d}x$?
	\choice
	{$\dfrac{2+\sqrt{2}}{2}$}
	{\True $\dfrac{2-\sqrt{2}}{2}$}
	{$\dfrac{\sqrt{2}}{2}$}
	{$-\dfrac{\sqrt{2}}{2}$}
	\loigiai{
		Ta có $I=\displaystyle \int\limits_{0}^{\tfrac{\pi}{4}} \sin x \mathrm{\,d}x=-\cos x \bigg|_0^{\tfrac{\pi}{4}} = - \cos \left(\dfrac{\pi}{4}\right)+\cos 0= -\dfrac{\sqrt{2}}{2}+1=\dfrac{2-\sqrt{2}}{2}$.
	}
\end{ex}
\begin{ex}%[Cau-8]%[2D4N2-2]
	Tính tích phân $I=\displaystyle \int\limits_{1}^{2}\dfrac{x^2+4x}{x}\mathrm{\,d}x$?
	\choice
	{$I=\dfrac{29}{2}$}
	{$I=-\dfrac{11}{2}$}
	{\True $I=\dfrac{11}{2}$}
	{$I=-\dfrac{29}{2}$}
	\loigiai{
		Ta có $I=\displaystyle \int\limits_{1}^{2} \dfrac{x^2+4x}{x}\mathrm{\,d}x= \int\limits_{1}^{2}(x+4)\mathrm{d}x=\dfrac{11}{2}$.
	}
\end{ex}
\begin{ex}%[Cau-9]%[2D4N2-3]
	Cho tích phân $I=\displaystyle \int\limits_{0}^{\tfrac{\pi}{2}}(4x-1+\cos x)\mathrm{d}x =\pi\left(\dfrac{\pi}{a} -\dfrac{1}{b}\right)+c$, $(a$, $b$, $c\in \mathbb{Q})$. Tính $a-b+c$.
	\choice
	{\True $1$}
	{$-2$}
	{$\dfrac{1}{3}$}
	{$\dfrac{1}{2}$}
	\loigiai{
		Ta có
		\begin{align*}
			I &= \displaystyle \int\limits_{0}^{\tfrac{\pi}{2}} \left(4x - 1 + \cos x\right) \mathrm{d}x \\
			&= \left(2x^2 - x + \sin x\right)\Big|_0^{\tfrac{\pi}{2}} \\
			&= 2\cdot \left(\dfrac{\pi}{2}\right)^2 - \dfrac{\pi}{2} + \sin \left(\dfrac{\pi}{2}\right) \\
			&= \dfrac{\pi^2}{2} - \dfrac{\pi}{2} + 1 \\
			&= \pi\left(\dfrac{\pi}{2}-\dfrac{1}{2}\right)+1.
		\end{align*}
		Suy ra $a=2$, $b=2$, $c=1$. \\
		Vậy $a-b+c=2-2+1=1$.
	}
\end{ex}
\begin{ex}%[Cau-10]%[2D4H2-2]
	Cho $I=\displaystyle \int\limits_{0}^{1}\left(4x-2m^2\right)\mathrm{d}x$. Có bao nhiêu giá trị nguyên của $m$ để $I+6>0$?
	\choice
	{$1$}
	{$5$}
	{$2$}
	{\True $3$}
	\loigiai{
		Ta có $I=\displaystyle \int\limits_{0}^{1}\left(4x-2m^2\right)\mathrm{d}x= \left(2x^2-2m^2 x\right)\bigg|_0^1=-2m^2+2$.\\
		Khi đó $I+6>0\Leftrightarrow-2m^2+2+6>0\Leftrightarrow-m^2+4>0\Leftrightarrow-2<m<2$.\\
		Mà $m$ là số nguyên nên $m \in\{-1;0;1\}$.\\
		Vậy có $3$ giá trị nguyên của $m$ thỏa mãn yêu cầu.
	}
\end{ex}
\begin{ex}%[Cau-11]%[2D4H2-1]
	Cho $\displaystyle \int\limits_{1}^{2} \left[3f(x) + 2g(x)\right]\mathrm{d}x=1$, $\displaystyle \int\limits_{1}^{2} \left[2f(x) - g(x)\right]\mathrm{d}x=-3$. Khi đó $\displaystyle \int\limits_{1}^{2} f(x) \mathrm{\,d}x$ bằng
	\choice
	{$\dfrac{11}{7}$}
	{\True $-\dfrac{5}{7}$}
	{$\dfrac{6}{7}$}
	{$\dfrac{16}{7}$}
	\loigiai{
		Ta có $\displaystyle \int\limits_{1}^{2} \left[3f(x) + 2g(x)\right] \mathrm{d}x = 1 \Leftrightarrow \displaystyle 3\int\limits_{1}^{2} f(x) \mathrm{\,d}x + \displaystyle 2\int\limits_{1}^{2} g(x) \mathrm{\,d}x= 1$.\\
		Và $\displaystyle \int\limits_{1}^{2} \left[2f(x) - g(x)\right] \mathrm{d}x = -3 \Leftrightarrow \displaystyle 2\int\limits_{1}^{2} 2f(x) \mathrm{\,d}x - \displaystyle \int\limits_{1}^{2} g(x) \mathrm{\,d}x = -3$.\\
		Đặt $a = \displaystyle \int\limits_{1}^{2} f(x) \mathrm{\,d}x$, $b = \displaystyle \int\limits_{1}^{2} g(x) \mathrm{\,d}x$ ta có hệ phương trình
		$$\heva{&3a+2b=1\\&2a-b=-3\\} \Leftrightarrow \heva{&a=-\dfrac{5}{7}\\&b=\dfrac{11}{7}.}$$
		Vậy $\displaystyle \int\limits_{1}^{2} f(x) \mathrm{\,d}x=\dfrac{5}{7}$.
	}
\end{ex}
\begin{ex}%[Cau-12]%[2D4V2-6]
	Một vật chuyển động chậm với vận tốc $v(t)=160-10t$ (m/s). Quãng đường mà vật di chuyển được từ thời điểm $t=0$ (s) đến thời điểm mà vật dừng lại là
	\choice
	{$160$ (m)}
	{\True $1280$ (m)}
	{$0$ (m)}
	{$144$ (m)}
	\loigiai{
		Vật dừng lại đồng nghĩa với $v(t)=0 \Leftrightarrow 160-10t=0 \Leftrightarrow t=16$ (s).\\
		Quãng đường vật đi được là $s(t)=\displaystyle \int\limits_{0}^{16}(160-10t) \mathrm{d}t = \left(160t-5t^2\right)\bigg|_0^{16}=1280$ (m).
	}
\end{ex}
\Closesolutionfile{ans}
\indapan{6}{ans/ans-2-B12-De2-NLC}
\Opensolutionfile{ans}[ans/ans-2-B12-De2-DS]
\TNTF
\begin{ex}%[Cau-1]%[2D4H2-1]
	Cho hàm số $y=f(x)$ liên tục trên $\mathbb{R}$ và thỏa mãn $\displaystyle \int\limits_{-1}^{10} f(x)\mathrm{\,d}x=15$,\break $\displaystyle \int\limits_{3}^{5} f(x) \mathrm{\,d}x=-2$, $\displaystyle \int\limits_{-1}^{12} f(x) \mathrm{\,d}x=5$.
	\choiceTF[t]
	{$\displaystyle \int\limits_{10}^{-1} 2f(x)\mathrm{\,d}x=30$}
	{$\displaystyle \int\limits_{10}^{12} \left[f(x)-2\right] \mathrm{d}x=-12$}
	{\True $\displaystyle \int\limits_{-1}^{3} f(x) \mathrm{\,d}x + \int\limits_{5}^{10} f(x) \mathrm{\,d}x=17$}
	{Biết rằng $f(x)>0$, $\forall x>3$; $f(x)<0$, $\forall x<3$ và $\displaystyle \int\limits_{-1}^{12} \left|f(x)\right|\mathrm{d}x=5$. Khi đó $\displaystyle \int\limits_{-1}^{3} f(x)\mathrm{\,d}x - \int\limits_{5}^{12} f(x)\mathrm{\,d}x=3$}
	\loigiai{
		
		\begin{itemchoice}
			\itemch Sai. Vì $ \displaystyle \int\limits_{10}^{-1} 2f(x) \mathrm{\,d}x = -2\int\limits_{-1}^{10} f(x) \mathrm{\,d}x = -2\cdot 15=-30$.
			\itemch Sai. Vì 
			\begin{align*}
				\displaystyle \int\limits_{10}^{12} \left[f(x)-2\right] \mathrm{d}x &= \int\limits_{10}^{12} f(x) \mathrm{\,d}x - \int\limits_{10}^{12} 2 \mathrm{\,d}x \\
				&= \int\limits_{-1}^{12} f(x) \mathrm{\,d}x - \int\limits_{-1}^{10} f(x) \mathrm{\,d}x - 4 \\
				&= 5-15-4 \\
				&= -14.
			\end{align*}
			\itemch Đúng. Vì $\displaystyle \int\limits_{-1}^{3} f(x) \mathrm{\,d}x + \int\limits_{5}^{10} f(x) \mathrm{\,d}x = \int\limits_{-1}^{10} f(x) \mathrm{\,d}x - \int\limits_{3}^{5} f(x) \mathrm{\,d}x = 15 - (-2) = 17$.
			\itemch Sai.\\
			Ta có $f(x)>0$, $\forall x>3$; $f(x)<0$, $\forall x<3$.\\
			Suy ra $\displaystyle \int\limits_{-1}^{12} \left|f(x)\right|\mathrm{d}x= -\int\limits_{-1}^{3} f(x)\mathrm{\,d}x+\int\limits_{3}^{5} f(x)\mathrm{\,d}x + \int\limits_{5}^{12} f(x)\mathrm{\,d}x=5$.\\
			Khi đó\\
			\begin{align*}
				\displaystyle \int\limits_{-1}^{3} f(x) \mathrm{\,d}x -  \int\limits_{5}^{12} f(x) \mathrm{\,d}x &=  \int\limits_{3}^{5} f(x) \mathrm{\,d}x - 5\\
				&= -2-5 \\
				&= -7.
			\end{align*}
		\end{itemchoice}
	}
\end{ex}
\begin{ex}%[Cau-2]%[2D4H2-2]
	Cho hàm số $y = f(x)$ liên tục trên $\mathbb{R}$ thỏa mãn $f(x) = \heva{&\dfrac{4x^2-3}{x}& \text{khi} &x\ge 1\\&ax+b &\text{khi} &-2<x<1\\&x^2+4x-4 &\text{khi} &x \le -2.}$
	\choiceTF[t]
	{\True $\displaystyle \int\limits_{-5}^{-2} f(x) \mathrm{\,d}x = -15$}
	{\True $\displaystyle \int\limits_{3}^{4} f(x) \mathrm{\,d}x = 14 +3\ln 3 - 6\ln 2$}
	{$\displaystyle \int\limits_{0}^{1} f(x) \mathrm{\,d}x = a+b$}
	{\True $\displaystyle \int\limits_{-3}^{0} f(x) \mathrm{\,d}x = \dfrac{-53}{3}$}
	\loigiai{
		\begin{itemchoice}
			\itemch Đúng. Vì $\displaystyle \int\limits_{-5}^{-2} f(x) \mathrm{\,d}x =\int\limits_{-5}^{-2} \left(x^2+4x-4\right) \mathrm{d}x = -15$.
			\itemch Đúng. Vì $\displaystyle \int\limits_{3}^{4} f(x) \mathrm{\,d}x = \int\limits_{3}^{4} \dfrac{4x^2-3}{x} \mathrm{\,d}x = \int\limits_{3}^{4} \left(4x - \dfrac{3}{x} \right) \mathrm{d}x = 14 +3\ln 3 - 6\ln 2$.
			\itemch Sai. Vì $\displaystyle \int\limits_{0}^{1} f(x) \mathrm{\,d}x =\int\limits_{0}^{1}(ax+b)\mathrm{d}x=\left(\dfrac{ax^2}{2} + bx\right)\bigg|_0^1=\dfrac{a}{2}+b\ne a+b$.\\ (Chỉ đúng với $a=0$).
			\itemch Đúng. Vì $\displaystyle \int\limits_{-3}^{0} f(x) \mathrm{\,d}x= \int\limits_{-3}^{-2} \left(x^2+4x-4\right)\mathrm{d}x+\int\limits_{-2}^{0} (ax+b)\mathrm{d}x$.\\
			Do hàm số $y=f(x)$ liên tục trên $\mathbb{R}$ nên 
			$$\heva{&\lim\limits_{x\to 1^-} f(x) = \lim\limits_{x\to 1^+} f(x)\\&\lim\limits_{x\to -2^+} f(x) = \lim\limits_{x\to -2^-} f(x)\\}\Rightarrow \heva{&a+b=1\\&-2a+b=-8\\} \Rightarrow a=3;\, b=-2.$$
			Suy ra $\displaystyle \int\limits_{-3}^{0} f(x) \mathrm{\,d}x = \int\limits_{-3}^{-2} \left(x^2+4x-4\right)\mathrm{d}x+\int\limits_{-2}^{0} (3x-2)\mathrm{d}x=\dfrac{-53}{3}$.
		\end{itemchoice}
	}
\end{ex}
\begin{ex}%[Cau-3]%[2D4H2-4]
	Cho hàm số $f(x)$; $g(x)$ thỏa mãn $\displaystyle \int\limits_{2}^{6} f(x) \mathrm{\,d}x= 3$; $\displaystyle \int\limits_{2}^{6} g(x)\mathrm{\,d}x=-2$.
	\choiceTF[t]
	{\True $\displaystyle \int\limits_{2}^{6} \left[f(x)+g(x)\right]\mathrm{,d}x=1$}
	{$\displaystyle \int\limits_{2}^{6} \left[3f(x)-g(x)-3\right]\mathrm{d}x=10$}
	{\True $\displaystyle \int\limits_{2}^{6} \left[3\mathrm{e}^x-2f(x)\right]\mathrm{d}x= 3\mathrm{e}^6-3\mathrm{e}^2-6$}
	{Biết $\displaystyle \int\limits_{2}^{6}\left[3g(x)-\dfrac{2x-3}{x^2}\right]\mathrm{d}x =a+b\ln 3$, với $a$; $b\in\mathbb{Q}$.
		Khi đó $a^2+12b=-8$}
	\loigiai{
		\begin{itemchoice}
			\itemch Đúng. Vì $\displaystyle \int\limits_{2}^{6}\left[f(x)+g(x)\right] \mathrm{d}x=\displaystyle\int\limits_{2}^{6} f(x)\mathrm{\,d}x+\displaystyle\int\limits_{2}^{6} g(x) \mathrm{\,d}x=3-2=1$.
			\itemch Sai. Vì 
			\begin{align*}
				&\displaystyle \int\limits_{2}^{6}\left[3f(x)-g(x)-3\right]\mathrm{d}x \\
				= & 3\displaystyle\int\limits_{2}^{6} f(x) \mathrm{\,d}x-\displaystyle\int\limits_{2}^{6} g(x) \mathrm{\,d}x - \displaystyle\int\limits_{2}^{6} 3 \mathrm{\,d}x \\
				= & 3\cdot 3-(-2)-12=-1\ne 10.
			\end{align*}
			\itemch Đúng. Vì $\displaystyle \int\limits_{2}^{6} \left[3\mathrm{e}^x-2f(x)\right] \mathrm{\,d}x = 3\displaystyle\int\limits_{2}^{6} \mathrm{e}^x \mathrm{\,d}x - 2\displaystyle\int\limits_{2}^{6} f(x) \mathrm{\,d}x = 3\mathrm{e}^6-3\mathrm{e}^2-6$.
			\itemch Sai. Ta có
			\begin{align*}
				& \displaystyle \int\limits_{2}^{6} \left[3g(x)-\dfrac{2x-3}{x^2}\right] \mathrm{d}x \\
				= &\ 3\displaystyle\int\limits_{2}^{6} g(x)\mathrm{\,d}x- \displaystyle\int\limits_{2}^{6} \dfrac{2x-3}{x^2} \mathrm{\,d}x=-6- \left(2\ln \left|x\right|+\dfrac{3}{x}\right)\bigg|_2^6 \\
				= &-5 + 2\ln 3.
			\end{align*}
			Suy ra $a=-5$; $b=2$.\\
			Vậy $a^2+12b=25+24=49\ne 25$.
		\end{itemchoice}
	}
\end{ex}
\begin{ex}%[Cau-4]%[2D4V2-2]
	Cho hàm số $y=f(x)$ liên tục trên $\mathbb{R}$, đồ thị hàm số $(C)\colon y=f'(x)$ trên đoạn $[-3;6]$ là đường gấp khúc như hình vẽ. Khi đó
	\begin{center}
		\begin{tikzpicture}[scale=0.6,>=stealth, font=\footnotesize, line join=round, line cap=round]  
			%Vẽ hệ trục
			\draw[->] (-4.5,0)--(7.5,0) node[below]{$x$};
			\draw[->] (0,-3.5)--(0,4.5) node[left]{$y$};
			\node at (0,0) [below left]{$O$};
			\path 
				(-3,-2) coordinate (A)
				(-3,0) coordinate (Ax)
				(0,-2) coordinate (Ay)
				(2,3) coordinate (B)
				(2,0) coordinate (Bx)
				(0,3) coordinate (By)
				(6,-1) coordinate (C)
				(6,0) coordinate (Cx)
				(0,-1) coordinate (Cy)
				(5,0) coordinate (Ex)
			;
			\fill[black] (A) circle(1pt) node[below left]{$A$};
			\fill[black] (Ax) circle(1pt) node[above]{$-3$};
			\fill[black] (Ay) circle(1pt) node[right]{$-2$};
			\fill[black] (B) circle(1pt) node[above right]{$B$};
			\fill[black] (Bx) circle(1pt) node[below]{$2$};
			\fill[black] (By) circle(1pt) node[left]{$3$};
			\fill[black] (C) circle(1pt) node[below]{$C$};
			\fill[black] (Cx) circle(1pt) node[above]{$6$};
			\fill[black] (Ex) circle(1pt) node[above right]{$5$};
			\fill[black] (Ex) circle(0pt) node[below left]{$E$};
			\draw[dashed] (Ax)--(A)--(Ay) (Bx)--(B)--(By) (Cx)--(C);
			\draw (A)--(B)--(C);
		\end{tikzpicture}
	\end{center}
	\choiceTF[t]
	{\True $\displaystyle \int\limits_{-3}^{-1} f'(x)\mathrm{\,d}x=-2$}
	{\True $\displaystyle \int\limits_{0}^{1} f'(x)\mathrm{\,d}x=\dfrac{3}{2}$}
	{$f(2)-f(6)=4$}
	{$f(5)+f(-3)-2f(2)=-10$}
	\loigiai{
		\begin{itemchoice}
			\itemch Đúng.\\
			Ta có $A(-3;-2)$, $B(2;3)$ $\Rightarrow AB\colon y=x+1$.\\
			Khi đó $\displaystyle \int\limits_{-3}^{-1} f'(x)\mathrm{\,d}x= \displaystyle\int\limits_{-3}^{-1} (x+1)\mathrm{\,d}x=-2$.
			\itemch Đúng. Vì $\displaystyle \int\limits_{0}^{1} f'(x) \mathrm{\,d}x= \displaystyle\int\limits_{0}^{1}(x+1)\mathrm{\,d}x=\dfrac{3}{2}$.
			\itemch Sai.\\
			Ta có $B(2;3)$, $E(5;0)\Rightarrow BC\colon y=-x+5$.\\
			Khi đó $\displaystyle \int\limits_{2}^{6} f'(x)\mathrm{\,d}x= \displaystyle\int\limits_{2}^{6}(-x+5) \mathrm{\,d}x = 4$.\\
			Vì vậy $f(2)-f(6)=-\displaystyle \int\limits_{2}^{6} f'(x)\mathrm{\,d}x=-4$.
			\itemch Sai.\\
			Ta có $\displaystyle \int\limits_{-3}^{2} f'(x)\mathrm{\,d}x= \displaystyle\int\limits_{-3}^{2} (x+1)\mathrm{\,d}x=\dfrac{5}{2}=f(2)-f(-3)$.\\
			Mặt khác $\displaystyle \int\limits_{2}^{5} f'(x)\mathrm{\,d}x= \displaystyle\int\limits_{2}^{5}(-x+5)\mathrm{\,d}x=\dfrac{9}{2}=f(5)-f(2)$.\\
			Vì vậy $f(5)+f(-3)-2f(2)=\dfrac{9}{2}-\dfrac{5}{2}=2$.
		\end{itemchoice}
	}
\end{ex}
\Closesolutionfile{ans}
\indapan{3}{ans/ans-2-B12-De2-DS}
\Opensolutionfile{ans}[ans/ans-2-B12-De2-KQ]
\TNSA
\begin{ex}%[Cau-1]%[2D4V2-2]
	Cho $\displaystyle \int\limits_{1}^{4}\sqrt{\dfrac{1}{4x}+\dfrac{\sqrt{x}+ \mathrm{e}^x}{\sqrt{x}\cdot \mathrm{e}^{2x}}}\mathrm{\,d}x=a+\mathrm{e}^b-\mathrm{e}^c$ với $a$, $b$, $c$ là các số nguyên. Tính giá trị của biểu thức $S=a+b+c$.
	\shortans{$-4$}
	\loigiai{
		Ta có
		\begin{align*}
			\displaystyle \int\limits_{1}^{4} \sqrt{\dfrac{1}{4x} + \dfrac{\sqrt{x} + \mathrm{e}^x}{\sqrt{x}\cdot \mathrm{e}^{2x}}} \mathrm{\,d}x 
			&= \displaystyle\int\limits_{1}^{4} \sqrt{\left(\dfrac{1}{2\sqrt{x}}\right)^2 + 2\dfrac{1}{2\sqrt{x}\cdot \mathrm{e}^x}+\left(\dfrac{1}{\mathrm{e}^x}\right)^2} \mathrm{\,d}x\\
			&= \displaystyle\int\limits_{1}^{4}\sqrt{\left(\dfrac{1}{2\sqrt{x}} + \dfrac{1}{\mathrm{e}^x}\right)^2}\mathrm{d}x\\
			&=\displaystyle\int\limits_{1}^{4}\left(\dfrac{1}{2\sqrt{x}}+\dfrac{1}{\mathrm{e}^x}\right) \mathrm{\,d}x\\
			&=\left(\sqrt{x}-\mathrm{e}^{-x}\right)\bigg|_1^4\\
			&=1-\mathrm{e}^{-4}+\mathrm{e}^{-1}\\
			&=a+\mathrm{e}^b-\mathrm{e}^c.
		\end{align*}
		$\Rightarrow \heva{&a=1\\&b=-1\\&c=-4.}$\\
		Vậy $a+b+c=1+(-1)+(-4)=-4$.
	}
\end{ex}
\begin{ex}%[Cau-2]%[2D4V2-6]
	Tốc độ chuyển động của thang máy từ tầng $1$ lên tầng cao nhất theo thời gian $t$ (giây) được cho bởi công thức 
	$$v(t) = \heva{&t&\text{khi}&\ 0 \le t \le 2\\&2 &\text{khi}&\ 2 < t \le 20\\&12-0{,}5t &\text{khi}&\ 20 < t \le 24.}$$
	Tính vận tốc trung bình của thang máy.
	\shortans{$1{,}75$}
	\loigiai{
		Quãng đường chuyển động của thang máy là\\
		$s=\displaystyle \int\limits_{0}^{24} v(t)\mathrm{\,d}t = \displaystyle\int\limits_{0}^{2} v(t)\mathrm{\,d}t + \displaystyle\int\limits_{2}^{20} v(t)\mathrm{\,d}t + \displaystyle\int\limits_{20}^{24} v(t)\mathrm{\,d}t = \displaystyle\int\limits_{0}^{2} t\mathrm{\,d}t +\displaystyle\int\limits_{2}^{20} 2 \mathrm{\,d}t+\displaystyle\int\limits_{20}^{24} (12-0{,}5t)\mathrm{d}t = 42$.\\
		Tốc độ trung bình của thang máy là $v_{tb}=\dfrac{s}{t}=\dfrac{42}{24}=1{,}75$ (m/s).
	}
\end{ex}
\begin{ex}%[Cau-3]%[2D4H2-3]
	Biết $\displaystyle \int\limits_{0}^{\tfrac{\pi}{2}} \left(x-1+\sin 2x\right) \mathrm{d}x = \pi \left(\dfrac{\pi}{a} - \dfrac{1}{b}\right) + 1$, $(a$, $b\in \mathbb{Q})$. Tính $a+2b$.
	\shortans{$12$}
	\loigiai{
		Ta có
		\begin{align*}
			\displaystyle \int\limits_{0}^{\tfrac{\pi}{2}} (x-1+\sin 2x) \mathrm{d}x &= \left(\dfrac{1}{2}x^2-x-\dfrac{1}{2}\cos 2x\right)\bigg|_0^{\tfrac{\pi}{2}}\\
			&=\dfrac{1}{2}\cdot \left(\dfrac{\pi}{2}\right)^2 - \dfrac{\pi}{2} - \dfrac{1}{2} \cos \left(2\cdot \dfrac{\pi}{2}\right) + \dfrac{1}{2}\\
			&=\dfrac{\pi^2}{8}-\dfrac{\pi}{2}+1\\
			&=\pi\left(\dfrac{\pi}{8}-\dfrac{1}{2}\right)+1.
		\end{align*}
		Suy ra $a=8$; $b=2$.\\
		Vậy $a+2b=8+2\cdot 2=12$.
	}
\end{ex}
\begin{ex}%[Cau-4]%[2D4V2-3]
	Cho $M$, $N$ là các số thực, xét hàm số $f(x)=M\cdot \sin \pi x + N\cdot \cos \pi x$ thỏa mãn $f(1)=3$ và $\displaystyle \int\limits_{0}^{\tfrac{1}{2}} f(x) \mathrm{\,d}x= -\dfrac{1}{\pi}$. Tính $f'\left(\dfrac{1}{4}\right)$. (Kết quả làm tròn đến hàng phần mười).
	\shortans{$11{,}1$}
	\loigiai{
		Ta có $f(1)=3\Leftrightarrow M\cdot \sin \pi+N\cdot \cos \pi=3\Leftrightarrow N=-3$.\\
		Mặt khác \\
		\begin{align*}
			\displaystyle \int\limits_{0}^{\tfrac{1}{2}} f(x) \mathrm{\,d}x = -\dfrac{1}{\pi} &\Leftrightarrow \int\limits_{0}^{\tfrac{1}{2}} \left(M\cdot \sin \pi x + N \cdot \cos \pi x\right) \mathrm{d}x = -\dfrac{1}{\pi}\\
			&\Leftrightarrow \left(-\dfrac{M}{\pi}\cos \pi x - \dfrac{3}{\pi} \sin \pi x\right)\bigg|_0^{\tfrac{1}{2}} = -\dfrac{1}{\pi}\\
			&\Leftrightarrow -\dfrac{3}{\pi} + \dfrac{M}{\pi} = -\dfrac{1}{\pi}\\
			&\Leftrightarrow M = 2.
		\end{align*}
		Ta được $f(x) = 2\cdot \sin \pi x - 3\cdot \cos \pi x$ nên $f'(x) = 2\pi \cos \pi x + 3\pi \sin \pi x$.\\
		Vậy $f'\left(\dfrac{1}{4}\right) = \dfrac{5\pi \sqrt{2}}{2} \approx 11{,}1$.
	}
\end{ex}
\begin{ex}%[Cau-5]%[2D4V2-6]
	Ba Tí muốn làm cửa sắt được thiết kế như hình bên dưới. Vòm cổng có hình dạng là một Parabol. Giá $1$ m$^2$ cửa sắt là $660\,000$ đồng. Cửa sắt có giá (nghìn đồng) là bao nhiêu?
	\begin{center}
		\begin{tikzpicture}[smooth,samples=300,scale=1,>=stealth, font=\footnotesize]
			%\draw[->] (-3.0,0)--(4.0,0) node[below]{$x$};
			%\draw[->] (0,-1.8)--(0,1.5) node[right]{$y$};
			\path 
			(-2.5,0) coordinate (A)
			(2.5,0) coordinate (B)
			(0,0.5) coordinate (C)
			(2.5,-1.5) coordinate (E)
			(-2.5,-1.5) coordinate (D)
			(0,0) coordinate (O)
			;
			\draw[thick, magenta,domain=-2.5:2.5] plot(\x,{-0.08*(\x)^2+1/2});
			%\draw[thick, magenta,domain=-2.5:2.5] plot(\x,{-0.04*(\x)^2+2-0.2});
			\draw[dashed] (D)--($(D)+(0,-0.5)$)--($(E)+(0,-0.5)$)--(E);
			\draw[dashed] (C)--(3,0.5)--(3,-1.5)--(0,-1.5);
			%\draw[dashed] (A)--(B);
			\draw (A)--(D)--(E)--(B) (0,-1.5)--(0,0.5);
			\node[rotate=90] at (3.2,-0.8) {$2$ m};
			\node[rotate=90] at (2.7,-0.8) {$1{,}5$ m};
			\node[rotate=0] at (0.3,-2.2) {$5$ m};
		\end{tikzpicture}
	\end{center}
	\shortans{$6050$}
	\loigiai{
		\begin{center}
			\begin{tikzpicture}[smooth,samples=300,scale=1,>=stealth, font=\footnotesize]
				\draw[->] (-3.0,0)--(4.0,0) node[below]{$x$};
				\draw[->] (0,-1.8)--(0,1.5) node[right]{$y$};
				\path 
					(-2.5,0) coordinate (A)
					(2.5,0) coordinate (B)
					(0,0.5) coordinate (C)
					(2.5,-1.5) coordinate (E)
					(-2.5,-1.5) coordinate (D)
					(0,0) coordinate (O)
				;
				\draw[thick, magenta,domain=-2.5:2.5] plot(\x,{-0.08*(\x)^2+1/2});
				%\draw[thick, magenta,domain=-2.5:2.5] plot(\x,{-0.04*(\x)^2+2-0.2});
				\draw[dashed] (D)--($(D)+(0,-0.5)$)--($(E)+(0,-0.5)$)--(E);
				\draw[dashed] (C)--(3,0.5)--(3,-1.5)--(0,-1.5);
				\draw[dashed] (A)--(B);
				\draw (A)--(D)--(E)--(B);
				\node[rotate=90] at (3.2,-0.8) {$2$ m};
				\node[rotate=90] at (2.7,-0.8) {$1{,}5$ m};
				\node[rotate=0] at (0.3,-2.2) {$5$ m};
				\draw[fill=black] (O) circle(1pt) node[below right]{$O$};
				\draw[fill=black] (A) circle(1pt) node[below left]{$A$};
				\draw[fill=black] (B) circle(1pt) ($(B)+(60:0.3)$) node {$B$};
				\draw[fill=black] (C) circle(1pt) node[above right]{$C$};
				\draw[fill=black] (D) circle(1pt) node[left]{$D$};
				\draw[fill=black] (E) circle(1pt) node[below right]{$E$};
				\draw[->] (-3.5,0.25) node[left]{Phần 1}--(-1,0.25);
				\draw[->] (-3.5,-1) node[left]{Phần 2}--(-2,-1);
			\end{tikzpicture}
		\end{center}
		Từ hình vẽ ta chia cửa rào sắt ra thành $2$ phần như trên.\\
		Khi đó $S = S_1 + S_2 = S_1 + 5\cdot 1{,}5 = S_1 + 7{,}5$.\\
		Để tính $S_1$ ta vận dụng kiến thức tính diện tích hình phẳng của tích phân.\\
		Gắn hệ trục $Oxy$ trong đó $O$ trung với trung điểm của $AB$, $OB\subset Ox$, $OC \subset Oy$.\\
		Theo đề bài ta có đường cong có dạng hình Parabol. Giả sử $(P)\colon y=ax^2+bx+c$.\\
		Khi đó $\heva{&A\left(-\dfrac{5}{2}; 0\right)\in (P)\\&B\left(\dfrac{5}{2}; 0\right)\in (P)\\&C\left(0;\dfrac{1}{2}\right)\in (P)} \Leftrightarrow\heva{&\dfrac{25}{4}a-\dfrac{5}{2}b+c=0\\&\dfrac{25}{4}a+ \dfrac{5}{2}b+c=0\\&c=\dfrac{1}{2}} \Leftrightarrow \heva{&a= -\dfrac{2}{25}\\ &b = 0\\&c=\dfrac{1}{2}.}$\\
		$\Rightarrow (P) \colon y = -\dfrac{2}{25} x^2 + \dfrac{1}{2}$.\\
		Diện tích $S_2 = 2\displaystyle\int\limits_{0}^{2{,}5} \left(-\dfrac{2}{25} x^2 + \dfrac{1}{2}\right) \mathrm{d}x = \dfrac{10}{6}\ (\mathrm{m}^2)$.\\
		$\Rightarrow S = \dfrac{55}{6}$ (m$^2$).\\
		Vậy  giá tiền cửa sắt là $\dfrac{55}{6} \cdot 660\,000= 6050$ (nghìn đồng).
	}
\end{ex}
\begin{ex}%[Cau-6]%[2D4V2-6]
	Một ô tô đang chạy đều với vận tốc $15$ (m/s) thì phía trước xuất hiện chướng ngại vật nên người lái đạp phanh gấp. Kể từ thời điểm đó, ô tô chuyển động chậm dần đều với gia tốc $-a$ (m/s$^2$). Tìm giá trị của $a$ biết ô tô chuyển động thêm được $20$ (m) thì dừng hẳn. (Kết quả làm tròn đến hàng phần trăm).
	\shortans{$5{,}63$}
	\loigiai{
		Gọi $x(t)$ là hàm biểu diễn quãng đường, $v(t)$ là hàm vận tốc.\\
		Ta có $\displaystyle \int\limits_{0}^{t} (-a) \mathrm{\,d}x = -at \Rightarrow v(t) = -at + 15$.\\
		Mặt khác $x(t)-x(0)=\displaystyle \int\limits_{0}^{t}v(t)\mathrm{\,d}x =\displaystyle \int\limits_{0}^{t}(-at+15)\mathrm{d}x=-\dfrac{1}{2}at^2 + 15t$.\\
		$\Rightarrow x(t)=-\dfrac{1}{2}at^2 + 15t$.\\
		Ta có $\heva{&v(t)=0\\&x(t)=20}\Leftrightarrow\heva{&-at+15=0\\&-\dfrac{1}{2}at^2 + 15t = 20}\Rightarrow -\dfrac{15}{2}t+15t=20 \Rightarrow t=\dfrac{8}{3}$.\\
		$\Rightarrow a=\dfrac{45}{8}\approx 5{,}63$.
	}
\end{ex}
\Closesolutionfile{ans}
\indapan{6}{ans/ans-2-B12-De2-KQ}